\chapter{Experimental Setup}
\label{ch:setup}

This chapter describes the data, the models and hardware, the configurations compared, and the metrics. It also
states the limits of the current setup, which apply to every result in Chapter~\ref{ch:results}.

\section{Dataset}

\subsection{Graphs and problems from the prior work}

The experiments so far reuse the dataset of the prior work \cite{skianis2024pseudocode}, generated by
\code{scripts/generate_dataset.py} and stored in \code{data/graphs/} and \code{data/graphs_questions/}. It contains
Erd\H{o}s--R\'enyi graphs in three sizes, 100 graphs per size (300 in total): small graphs with 5--10 nodes, medium
graphs with 11--20 nodes and large graphs with 21--50 nodes. For each graph the generator draws the edge probability
uniformly from $[0, 1]$, removes isolated nodes and self-loops, and redraws until the graph has the intended number of
nodes. All graphs are undirected and unweighted, and five small graphs are disconnected, so a ``minimum spanning
tree'' there is any spanning forest.

The ten problems are node count, edge count, node degree, neighbors of a node (\code{connected_nodes}), edge
existence, connectivity between two nodes, number of connected components, cycle check, shortest path, and minimum
spanning tree. The harness keeps only the parameters of each stored question (node identifiers) and rebuilds the
question from its own template (Table~\ref{tab:templates} in Appendix~\ref{app:tools}), because the stored questions
contain the full edge list. Topological sort and bipartiteness have tools but no questions yet; they come with the
scaled generator of M6.

\subsection{An audit of the dataset}
\label{sec:audit}

Counting on the development graphs showed that several tasks have, in practice, one possible answer
(Table~\ref{tab:audit}). The large graphs are very dense, between 63 and 1,005 edges on 21--50 nodes, so paths are
short (the average shortest-path length in the large questions is 1.5, at most 3) and almost every pair of nodes is
connected by an edge. A model that always answers ``yes'' or ``1'' would score about 100\% on edge existence, cycle
check and component count on the large graphs. These tasks still test whether the loop and tools work, but they
cannot separate a good model from a lucky one. Only connectivity on small graphs is balanced, and there are no
connectivity questions for large graphs at all.

The large graphs are a real test in three places: \code{connected_nodes} (on average 21 and up to 45 neighbors to
list), \code{node_degree} (counting up to 45), and \code{mst} (on average 36 and up to 49 edges to list). Those are the
tasks where copying mistakes, and the effect of verifiers, can show.

\begin{table}[ht]
  \centering
  \small
  \caption{Answer balance on the development graphs (0--49).}
  \label{tab:audit}
  % source: docs/learning_log.md (2026-10-06, dataset audit)
  \begin{tabular}{lrr}
    \toprule
    & small (5--10 nodes) & large (21--50 nodes) \\
    \midrule
    \code{edge_existence} answer is ``yes'' & 97 / 100 & 100 / 100 \\
    graph has a cycle & 47 / 50 & 50 / 50 \\
    graph is connected (one component) & 47 / 50 & 50 / 50 \\
    connectivity questions (yes / no) & 28 (14 / 14) & 0 \\
    shortest-path length & -- & average 1.5, maximum 3 \\
    \bottomrule
  \end{tabular}
\end{table}

Two consequences follow for the thesis: per-task accuracy should be reported next to a majority-answer baseline, so
that 100\% on a one-answer task is read correctly, and the M6 generator should balance yes/no answers and control
density, so that paths are long and not every pair is an edge.
\todo[inline]{Add the majority-answer baseline column to the report (planned, not implemented yet).}

\subsection{Generated tasks}

Four tasks have no questions in the dataset; their questions are generated from the development graphs with a fixed
seed per task, size and graph, so that every run sees the same questions.

\paragraph{Waypoint route (\texttt{shortest\_path\_via}).} ``What is the shortest route from node A to node B if it must
pass through node C?'' No single tool answers it: the ideal run calls \code{shortest_path(A, C)} and
\code{shortest_path(C, B)} and joins the two paths. The generator only keeps waypoints that force a detour, so the
shortcut \code{shortest_path(A, B)} is always wrong, and sometimes the best route revisits a node. It is fully
verified.

\paragraph{The combine family.} Three tasks that need one lookup to find candidates and then many lookups and a
comparison:
\begin{itemize}
  \item \code{hop_max_degree}: among the nodes within $k$ hops of a node, which has the highest degree (ties to the
        smallest identifier); $k = 1$ on the dense large graphs, where two hops already reach almost every node;
  \item \code{common_neighbors_max}: which other node shares the most neighbors with a given node;
  \item \code{triangle_count}: how many triangles include a given node.
\end{itemize}
The questions pick a node with at least three neighbors, so there is something to loop over. With tools alone these
tasks need dozens of calls; one loop in code over the tool functions does it in one step. Their answers are numbers,
so they are unverified by design.

\paragraph{Multi-step probes.} Two hand-written questions (``which node is farthest from node 0?'' and ``how many other
nodes are within 2 hops of node 5?'') were used once to test the \code{distances_from} and \code{neighborhood} tools.

\section{Development and test split}

Graphs 0--49 of each size are the development split and graphs 50--99 the test split. Tools, prompts and routing are
designed on the development split only. Every result in this draft is on the development split; the test split has
not been used.

\section{Models and hardware}

\subsection{Local models}

All harness runs so far use local models served by Ollama and called through LiteLLM: \code{qwen3:8b} (the default
cheap model so far) and \code{qwen3.5:9b} (used for the combine experiments). The model's default temperature is used;
no temperature has been set.
\todo[inline]{Decide and fix the sampling temperature before the final experiments, and report it.}

Real runs are made on a machine with an NVIDIA RTX 5070 with 12~GB of memory, where a simple question takes about 4
seconds. A MacBook Air with an M4 chip and 24~GB of memory also runs qwen3:8b, but about 4.5 times slower (about 30
seconds per answer in the pilot), and is used only for quick tests. Which machine a result comes from is stated with
it. Since 8 October 2026, every Ollama call requests a 16,384-token context; earlier runs used Ollama's default of
4,096 tokens (Section~\ref{sec:case-4k}).

The choice of the cheap model is not final. A pilot of four candidates that fit the 12~GB GPU is planned
(Table~\ref{tab:candidates}), each on 54 questions (all nine large-graph tasks, six questions each, one run), to
decide on strict accuracy, rescued tool calls, verifier rejections, and tokens and time per question.

\begin{table}[ht]
  \centering
  \small
  \caption{Candidate models for the planned pilot. Sizes are from model pages and blogs and are unverified.}
  \label{tab:candidates}
  % source: CLAUDE.md ("Model candidates"), docs/learning_log.md (2026-10-08, model choice)
  \begin{tabularx}{\textwidth}{l>{\raggedright\arraybackslash}X>{\raggedright\arraybackslash}p{3.5cm}}
    \toprule
    Model & Why & Fits 12 GB? \\
    \midrule
    \code{qwen3.5:9b} & Likely new default: successor to qwen3:8b, same memory (6.6--7.6 GB) & yes \\
    \code{gemma4:12b} & A second model family, native function calling, 7.7--8 GB & yes \\
    \code{gpt-oss:20b} & Strong function calling; 14 GB download, sources disagree on fit & to test \\
    \code{qwen3:8b} & Current baseline, rerun with today's tool set & yes \\
    \code{qwen3.6:35b-a3b} & Strongest local option (mixture of experts, 3B active) for M5 escalation & offload only \\
    \bottomrule
  \end{tabularx}
\end{table}

Hosted models are planned as stronger baselines in our harness: a frontier closed model and a strong open model
through an API.
\todo[inline]{Fix the hosted models (CLAUDE.md lists Sonnet 5.5 and Kimi K2.6 via OpenRouter) and the API access; there is no API access yet.}

\subsection{The general-purpose harness}

The comparisons with a general-purpose harness use Claude Code in headless mode (\code{claude -p}) with a Sonnet
model on a subscription; the reported dollar cost is the API-equivalent cost that Claude Code reports. Two setups were
used. In the first, Claude Code may only run the project's Python (with networkx) and read the graph file, and answers
with JSON on its last line. In the second, all of Claude Code's built-in tools are switched off and it sees only our
tools through the MCP server; it is run outside the repository so that no project instructions are loaded. In one run
the alias \code{sonnet} resolved to a different model version than intended, so model identifiers must be pinned in
the M6 baselines.

\section{Configurations}

\paragraph{Lenient and strict.} Every run is lenient (text tool calls are rescued), and the strict accuracy is computed
from the same run by counting every answer that needed a rescue as wrong.

\paragraph{Verifiers on and off.} With \code{--no-verify} the model is still asked for the same answers with the same
evidence fields; only the checking is switched off.

\paragraph{Tool ablations.} \code{--without-tool} removes a tool from the default set (for example \code{has_edge} or
\code{degree}) and \code{--with-tool} adds an opt-in tool. Because the default tool set grew during the project,
results from different dates are not directly comparable; each experiment states its tool set or prompt version.

\paragraph{Code versus tools.} Five setups compare tool calls with code on the combine tasks (Table~\ref{tab:setups}).
Setup B offers \code{run_python} next to the tools; setup H adds the strategy hint \code{CODE_HINT} to the system
prompt (``use a direct tool for a single lookup; when the answer needs many lookups, write one \code{run_python} call
that loops over the tool functions''); setup C offers code only, with our tools as functions; setup NX offers code
only, with networkx.

\begin{table}[ht]
  \centering
  \small
  \caption{Setups of the code-versus-tools experiments.}
  \label{tab:setups}
  % source: results/harness_runs/combine16k_qwen35_*/traces.jsonl (run_start), docs/architecture.md
  \begin{tabularx}{\textwidth}{l>{\raggedright\arraybackslash}X l l}
    \toprule
    Setup & Offered & \code{python} & hint \\
    \midrule
    A & Graph tools only & off & no \\
    B & Graph tools and \code{run_python} (tool functions inside code) & \code{tools} & no \\
    H & As B, plus \code{CODE_HINT} in the system prompt & \code{tools} & yes \\
    C & \code{run_python} only, tool functions inside code, no networkx & \code{tools} & no \\
    NX & \code{run_python} only, with \code{G} and \code{nx} (networkx) & \code{networkx} & no \\
    \bottomrule
  \end{tabularx}
\end{table}

\section{Metrics}

\begin{itemize}
  \item \emph{Accuracy} (lenient): the share of all answers, over every run, graded correct by \code{eval/tasks.py}.
        The grader accepts any shortest path, neighbors in any order, any spanning forest, and never counts
        \code{true} as 1.
  \item \emph{95\% confidence interval}: bootstrap over questions, with all runs of a question resampled together.
  \item \emph{Strict accuracy}: accuracy without answers rescued from text tool calls.
  \item \emph{Checked}: the share of answers a verifier actually looked at; \emph{rejected}: answers a verifier sent
        back; \emph{no answer}: runs that ended without one (\code{cannot_answer}, loop detected, step limit, context
        full, error).
  \item \emph{Cost}: prompt plus completion tokens per answer, and wall-clock time per answer. The dollar cost of local
        runs is zero; for Claude Code the reported API-equivalent cost is used.
  \item \emph{Process metrics}: tool precision, recall, F1, parameter match and exact match
        (Section~\ref{sec:process-metrics}), as in the GDS Agent benchmark.
  \item \emph{Code use}: the share of questions with at least one \code{run_python} call, with both code and graph
        tools, code calls and graph-tool calls per question, the share of code calls that ended in an error, and
        empty replies.
\end{itemize}

The primary comparison of the thesis, the accuracy-versus-cost Pareto front, combines accuracy with tokens, time and
dollars. No Pareto analysis has been done yet; it needs the hosted-model baselines and repeated runs (M7).

\section{Limitations of the current setup}

The results in Chapter~\ref{ch:results} should be read with these limitations in mind:
\begin{itemize}
  \item Most experiments are a single run, many on 3--15 questions; the plan requires at least three runs per
        configuration. With 15 questions, one question is 6.7 percentage points.
  \item Confidence intervals are reported where the report computes them, but with one run and few questions they
        are wide, and for tasks at 100\% the bootstrap interval collapses to a point.
  \item All results are on the development split, which was also used to design the tools and prompts.
  \item The existing dataset is saturated on large graphs (Section~\ref{sec:audit}); it cannot yet separate setups on
        most tasks.
  \item The tool set and prompt version changed between experiments, so experiments from different dates are not
        directly comparable.
  \item The comparisons with Claude Code use a different model and, in one case, different tools; they are
        previews, not the fair comparison planned for M6.
  \item The temperature was left at the model default.
\end{itemize}
